\documentclass[11pt]{article}
\usepackage[margin=1in]{geometry}
\usepackage{amsmath,amssymb,amsthm,microtype}
\usepackage[hidelinks]{hyperref}

\newcommand{\R}{\mathbb R}
\newcommand{\Mp}{\mathcal M_p}
\newcommand{\D}{\Delta}

\title{The two-dimensional $L^p$-Mahler conjectures\\
are answered by arXiv:2401.10992}
\author{Literature-status note for arXiv:2304.14363}
\date{13 August 2026}

\begin{document}
\maketitle

\begin{center}
\fbox{\parbox{0.91\textwidth}{\textbf{Status.}
Explicit later-paper answer to the complete two-dimensional subcase.  The
higher-dimensional conjectures remain open and imply the classical Mahler
conjectures.  This note records a literature identification, not a new proof.}}
\end{center}

\section{The source conjectures}

Berndtsson, Mastrantonis, and Rubinstein define, for a convex body
$K\subset\R^n$ and $p>0$,
\[
 h_{p,K}(y)=\frac1p\log\left(\frac1{|K|}
 \int_K e^{p\langle x,y\rangle}\,dx\right),
 \qquad
 \Mp(K)=|K|\int_{\R^n}e^{-h_{p,K}(y)}\,dy.
\]
Their Conjectures 1.3 and 1.4, stated on PDF page 7 of
arXiv:2304.14363v1 \cite{BMR}, ask for the following lower bounds.

\begin{itemize}
\item For every symmetric convex body $K\subset\R^n$,
\begin{equation}\label{eq:sym}
 \Mp([-1,1]^n)\leq \Mp(K).
\end{equation}
\item For every convex body $K\subset\R^n$,
\begin{equation}\label{eq:nonsym}
 \inf_{x\in\D_n}\Mp(\D_n-x)\leq\Mp(K),
\end{equation}
where $\D_n$ is the standard simplex.
\end{itemize}

The source formulates these for $p\in(0,\infty]$, with $p=\infty$ recovering
the classical symmetric and nonsymmetric Mahler conjectures.  It also proves
that the family of finite-$p$ conjectures implies the classical conjectures
\cite[Lemma 3.12]{BMR}.

\section{The explicit later answer in dimension two}

Mastrantonis and Rubinstein's follow-up
\emph{Two-dimensional B\l{}ocki, $L^p$-Mahler, and Bourgain conjectures},
arXiv:2401.10992v3 \cite{MR}, restates \eqref{eq:sym} and
\eqref{eq:nonsym} as its Conjectures 1.1 and 1.2.  Its main results on PDF
page 3 are exactly:
\begin{quote}
\textbf{Theorem 1.3.} Conjecture 1.1 holds when $n=2$.

\textbf{Theorem 1.4.} Conjecture 1.2 holds when $n=2$.
\end{quote}
The paper treats every finite $p>0$, in both the symmetric and nonsymmetric
settings.  It also notes immediately afterward that the endpoint $p=\infty$
in dimension two is Mahler's classical theorem.  Consequently the entire
$n=2$ specialization of source Conjectures 1.3--1.4 is resolved for
$p\in(0,\infty]$.

The identification is direct rather than inferred: the follow-up cites the
source conjectures by number, reproduces their formulas, and labels its
Theorems 1.3--1.4 as their two-dimensional verification.  The proof uses
Mahler sliding of polygon vertices and controls the non-polyhedral
$L^p$-polar along the slide.

\section{What remains open}

This is a complete answer only in dimension two.  No bounded search through
13 August 2026 found a general higher-dimensional resolution.  Attempting the
unrestricted source target inside this lane would amount to attacking the
classical Mahler conjectures, since the source's finite-$p$ family implies
them.  The responsible disposition is therefore:
\[
 \boxed{\text{$n=2$: answered by arXiv:2401.10992;
 higher dimensions: major open problem.}}
\]

The local source corpus, exact title/conjecture searches, and later
$L^p$-Mahler papers were checked.  The 2024 two-dimensional paper is the
precise match; later functional and heat-flow work does not claim the full
higher-dimensional inequalities.

\begin{thebibliography}{9}
\bibitem{BMR}
B.~Berndtsson, V.~Mastrantonis, and Y.~A.~Rubinstein,
\emph{$L^p$-polarity, Mahler volumes, and the isotropic constant},
arXiv:2304.14363v1 (2023).

\bibitem{MR}
V.~Mastrantonis and Y.~A.~Rubinstein,
\emph{Two-dimensional B\l{}ocki, $L^p$-Mahler, and Bourgain conjectures},
arXiv:2401.10992v3; to appear in Indiana Univ. Math. J.
\end{thebibliography}

\end{document}
